\subsection{极集与正交子空间}

定义 2. --- 设 F 和 G 是两个成对偶的（实）向量空间。对 F 的每个集 M，我们称由那些 \(y \in G\) （它们满足 \(\langle x, y \rangle \geqslant -1\) ，对所有 \(x \in M\) ）组成的集为 M 的极集。（对复向量空间，参见 II, 第 64 页。）

若 \(G_{1}\) 、 \(G_{2}\) 是 \(F^{*}\) 的两个满足 \(G_{1} \subset G_{2}\) 的子空间，则 M 在 \(G_{1}\) 中的极集是 \(G_{1}\) 与 M 在 \(G_{2}\) 中的极集的交。

当没有混淆危险时，我们用 \(M^{\circ}\) 表示 F 的子集 M 在 G 中的极集。类似地，我们定义 G 中的集在 F 中的极集。

显然，对每个标量 \(\lambda \neq 0\) 及所有 \(\mathbf{M} \subset \mathbf{F}\) ，有 \((\lambda \mathbf{M})^{\circ} = \lambda^{-1} \mathbf{M}^{\circ}\) 。关系 \(\mathbf{M} \subset \mathbf{N} \subset \mathbf{F}\) 蕴含 \(\mathbf{N}^{\circ} \subset \mathbf{M}^{\circ}\) ；若 \(\mathbf{N}\) 吸收 \(\mathbf{M}\) ，则 \(\mathbf{M}^{\circ}\) 吸收 \(\mathbf{N}^{\circ}\) ；对每个族 \((\mathbf{M}_{\alpha})\) （\(\mathbf{F}\) 的集组成的），\(\bigcup_{\alpha} \mathbf{M}_{\alpha}\) 的极集是诸极集 \(\mathbf{M}_{\alpha}^{\circ}\) 的交。由于对 \(y \in \mathbf{M}^{\circ}\) ，由关系 \(\langle x, y \rangle \geqslant -1\) 定义的闭半空间包含 0 和 \(\mathbf{M}\) ，我们看出，若 \(\mathbf{M}_{1}\) 是 \(\mathbf{M} \cup \{0\}\) 的凸包，则 \(\mathbf{M}_{1}^{\circ} = \mathbf{M}^{\circ}\) 。

显然 \(\mathbf{M} \subset \mathbf{M}^{\circ \circ}\) 。于是

\[
(\mathbf {M} ^ {\circ \circ}) ^ {\circ} \subset \mathbf {M} ^ {\circ} \subset (\mathbf {M} ^ {\circ}) ^ {\circ \circ} = (\mathbf {M} ^ {\circ \circ}) ^ {\circ}
\]

即 \(\mathbf{M}^{\circ \circ \circ} = \mathbf{M}^{\circ}\) （参见 S, III, § 1.5, prop. 2）。

若 \(\mathbf{M}\) 是 \(\mathrm{F}\) 的一个对称子集，则 \(\mathbf{M}^{\circ}\) 是 \(\mathbf{G}\) 的一个对称子集；这时 \(\mathbf{M}^{\circ}\) 也是由那些 \(y \in \mathbf{G}\) （它们满足 \(|\langle x, y \rangle| \leqslant 1\) ，对所有 \(x \in \mathbf{M}\) ）组成的集。

命题 4. --- (i) 对 F 的任一集 M，极集 \(\mathbf{M}^{\circ}\) 是一个包含 0 的凸集，且在 G 中对拓扑 \(\sigma(\mathbf{G},\mathbf{F})\) 是闭的。

\begin{enumerate}
\def\labelenumi{(\roman{enumi})}
\setcounter{enumi}{1}
\item
  若 \(\mathbf{M}\) 是一个以 0 为顶点的锥，则 \(\mathbf{M}^{\circ}\) 是一个以 0 为顶点的锥，它也是由那些 \(y\in \mathbf{G}\) （它们满足 \(\langle x,y\rangle \geqslant 0\) ，对所有 \(x\in \mathbf{M}\) ）组成的集。
\item
  若 \(\mathbf{M}\) 是 \(\mathbf{F}\) 的一个向量子空间，则 \(\mathbf{M}^{\circ}\) 是 \(\mathbf{G}\) 的一个向量子空间，它也是由那些 \(y \in \mathbf{G}\) （它们满足 \(\langle x, y \rangle = 0\) ，对所有 \(x \in \mathbf{M}\) ）组成的集。
\item
  由于线性型 \(y \mapsto \langle x, y \rangle\) 对 \(\sigma(\mathbf{G}, \mathbf{F})\) 连续，结论由定义以及由超平面确定的半空间是凸的这一事实立即得出。
\item
  若 \(\mathbf{M}\) 是一个以 0 为顶点的锥，且 \(x \in \mathbf{M}\) 、\(y \in \mathbf{M}^{\circ}\) ，则由于 \(\lambda x \in \mathbf{M}\) （对所有 \(\lambda > 0\) ），我们有 \(\langle \lambda x, y \rangle \geqslant -1\) ，即 \(\lambda \langle x, y \rangle \geqslant -1\) 。由于这对所有 \(\lambda > 0\) 成立，推出 \(\langle x, y \rangle \geqslant 0\) ，(ii) 得证。
\item
  类似地，若 \(\mathbf{M}\) 是 \(\mathbf{F}\) 的一个向量子空间，则关系 \(x \in \mathbf{M}\) 、\(y \in \mathbf{M}^{\circ}\) 这一次蕴含 \(\lambda \langle x, y \rangle \geqslant -1\) （对所有实数 \(\lambda\) ），这只有当 \(\langle x, y \rangle = 0\) 时才可能。
\end{enumerate}

若 \(\mathbf{M}\) 是 \(\mathbf{F}\) 的一个向量子空间，我们称 \(\mathbf{M}^{\circ}\) 为 \(\mathbf{M}\) 在 \(\mathbf{G}\) 中的正交（子空间）；若 \(\mathbf{G} \subset \mathbf{F}^{*}\) ，则 \(\mathbf{M}^{\circ}\) 是 \(\mathbf{G}\) 与 \(\mathbf{M}\) 在代数对偶 \(\mathbf{F}^{*}\) （\(\mathbf{F}\) 的）中的正交子空间的交 (A, II, § 2.4, def. 4)。

对 F 的向量子空间 M 和 G 的向量子空间 N，我们说 M 和 N 正交，如果 \(M \subset N^{0}\) （或者等价地，\(N \subset M^{\circ}\) ）。

定理 1（双极定理）. --- 设 F, G 是两个成对偶的实向量空间。对 F 的每个子集 M，\(M^{\circ\circ}\) ——M 在 G 中的极集 \(M^{\circ}\) 在 F 中的极集—— 是（对 \(\sigma(F, G)\) 而言）\(M \cup \{0\}\) 的闭凸包。

我们已看到，只须考虑 M 是凸的且 \(0 \in M\) 的情形。M 在拓扑 \(\sigma(F, G)\) 中的闭包记作 \(\overline{M}\) ，则 \(\overline{M}\) 是 F 中的一个凸集；II, 第 44 页的 prop. 4 表明 \(M^{\circ\circ} \supset \overline{M}\) 。另一方面，若 \(a \in F\) 不属于 \(\overline{M}\) ，则存在 F 中的一个闭超平面 H 把 a 与 \(\overline{M}\) 严格分离 (II, 第 38 页, prop. 4)；由于 H 不包含 0，存在 \(y \in G\) 使 H 有方程 \(\langle x, y \rangle = -1\) (II, 第 43 页, prop. 3)；于是 \(\langle x, y \rangle > -1\) （对所有 \(x \in \overline{M}\) ）而 \(\langle a, y \rangle < -1\) 。这蕴含 \(y \in M^{\circ}\) 且 \(a \notin M^{\circ\circ}\) ，从而关系 \(M^{\circ\circ} = \overline{M}\) 得证。

推论 1. --- 对任一族 \((\mathbf{M}_{\alpha})\) （诸集为 F 中对拓扑 \(\sigma(\mathrm{F}, \mathrm{G})\) 闭的凸集，且每个都包含 0），交 \(M = \bigcap_{\alpha} M_{\alpha}\) 的极集是（对 \(\sigma(\mathrm{G}, \mathrm{F})\) 而言）诸 \(M_{\alpha}^{\circ}\) 的并的闭凸包。

因为，若 N 是这个闭凸包，则

\[
\mathbf {N} ^ {\circ} = \bigcap_ {\alpha} \mathbf {M} _ {\alpha} ^ {\circ \circ} = \bigcap_ {\alpha} \mathbf {M} _ {\alpha} = \mathbf {M}
\]

于是 \(N = N^{\circ\circ} = M^{\circ}\) 。

若诸 \(M_{\alpha}\) 不是凸的，cor. 1 的结论未必成立。

推论 2. --- 对 F 的每个向量子空间 M，子空间 \(M^{\circ\circ}\) 是 M 在拓扑 \(\sigma(\mathrm{F}, \mathrm{G})\) 中的闭包。

注记. --- 拓扑 \(\sigma(\mathrm{G},\mathrm{F})\) 下 G 中 0 的每个邻域包含一个由有限多个形如 \(|\langle x_i,y\rangle |\leqslant 1\) （ \(1\leqslant i\leqslant n\) ）的不等式定义的邻域 V，其中 \(x_{i}\) 是 F 的任意点。若 A 是诸 \(x_{i}\) 组成的集的对称凸包，则 V 是 A 在 G 中的极集 \(\mathrm{A}^{\circ}\) 。我们可以说，F 中有限对称集（或其凸包）在 G 中的极集构成 \(\sigma(\mathrm{G},\mathrm{F})\) 下 G 中 0 的一个基本邻域系。若对偶在 F 中分离，则这些凸包对 \(\sigma(\mathrm{F},\mathrm{G})\) 是紧的 (II, 第 14 页, prop. 15 的 cor. 1) 且是有限维的。反之，F 中每个有限维紧凸集 C（带 \(\sigma(\mathrm{F},\mathrm{G})\) 拓扑）含于 F 的一个有限子集的凸包中。因为，设 M 是包含 C 的一个有限维向量子空间。若 \((e_i)_{1\leqslant i\leqslant n}\) 是 M 的一个基，我们可以假设 C 含于以 0 为中心、在基 \(e_i\) 的向量上构造的闭平行多面体中 (GT, VI, § 1.3)；而立即可见，这个平行多面体是点 \(\sum_{i=1}^{n} \varepsilon_i e_i\) （\(\varepsilon_i = \pm 1\) ）的凸包。

于是我们可以说，（若 \(\sigma(\mathrm{F},\mathrm{G})\) 是 Hausdorff 的）F 中有限维凸紧集（对 \(\sigma(\mathrm{F},\mathrm{G})\) 或对 F 上任何比 \(\sigma(\mathrm{F},\mathrm{G})\) 细的 Hausdorff 局部凸拓扑）的极集构成 \(\sigma(\mathrm{G},\mathrm{F})\) 下 0 的一个基本邻域系。

推论 3. --- 设 T 是局部凸空间 E 的拓扑，E' 是它的对偶 (II, 第 42 页, def. 1)。

\begin{enumerate}
\def\labelenumi{(\roman{enumi})}
\item
  \(\mathbf{E}\) 中的闭凸集对拓扑 \(\mathcal{T}\) 和对弱拓扑 \(\sigma (\mathrm{E},\mathrm{E}^{\prime})\) 是相同的。
\item
  对每个子集 \(\mathbf{M}\) （\(\mathbf{E}\) 的），\(\mathbf{M}^{\circ \circ}\) （在 \(\mathbf{E}\) 中；\(\mathbf{M}^{\circ}\) —— \(\mathbf{M}\) 在 \(\mathbf{E}'\) 中的极集—— 的极集）是 \(\mathbf{M} \cup \{0\}\) 对拓扑 \(\mathcal{T}\) 的闭凸包。
\end{enumerate}

显然，(ii) 由 (i) 及 th. 1 得出。由对偶 E' 的定义，II, 第 43 页的 prop. 3 表明 E 上对拓扑 T 的连续线性型与对 \(\sigma(\mathrm{E}, \mathrm{E}')\) 的连续线性型相同。因此 E 中的闭半空间对 T 和对 \(\sigma(\mathrm{E}, \mathrm{E}')\) 是相同的 (II, 第 15 页, prop. 17)，于是断言 (i) 由 II, 第 38 页的 cor. 1 得出。

